% 教材查阅指南模板（只能带教材）。编译：python scripts/compile_tex.py <文件>.tex
% 用途：考场上只能翻教材，所以这份指南告诉你每个考点在教材的哪一页、找什么。不抄教材原文。
% 页码一律用教材上印的页码（书页码），不是PDF页码；没有书页码时写章节号和小节标题，并在说明里提醒用户核对
% 规则：只改标题信息和正文，导言区（\begin{document}之前）不要动，不要自己加宏包和宏，不引入外部图片
% 【 】里是要替换的占位内容（替换时连括号一起去掉）。公式外的 % & # _ $ ^ ~ { } \ 要转义，见 references/exam-notes.md
\documentclass[a4paper,fontset=windows]{ctexart}
\ifdefined\SBSIZE\else\def\SBSIZE{10}\fi
\ifdefined\SBLEAD\else\def\SBLEAD{14}\fi
\ifdefined\SBMARGIN\else\def\SBMARGIN{18}\fi
\usepackage[margin=\SBMARGIN mm]{geometry}
\usepackage{amsmath,amssymb,enumitem,booktabs,longtable,xcolor}
\makeatletter
\renewcommand\normalsize{\@setfontsize\normalsize{\SBSIZE}{\SBLEAD}%
  \abovedisplayskip 4\p@ \@plus2\p@ \@minus2\p@
  \belowdisplayskip \abovedisplayskip
  \abovedisplayshortskip 2\p@ \@plus1\p@
  \belowdisplayshortskip 2\p@ \@plus1\p@
  \let\@listi\@listI}
\makeatother
\normalsize
\pagestyle{plain}
\setlength{\parindent}{0pt}
\setlength{\parskip}{3pt}
\setlist{nosep,leftmargin=1.4em}
\newlength{\sbw}
\newcommand{\kw}[1]{\ifdefined\SBFILL\settowidth{\sbw}{#1}\underline{\hspace*{\dimexpr\sbw+0.6em\relax}}\else\textbf{#1}\fi}

\begin{document}
\begin{center}
{\LARGE\bfseries 【课程名】　教材查阅指南}\\[4pt]
{\small 教材：【书名，版次】　考试：【YYYY-MM-DD】　规则：只能携带教材}
\end{center}

页码均为教材上印的页码。【如果是识别不到书页码的情况，写：本表按章节号和小节标题定位，页码需要对照教材核对。】

% ---- 按考点权重从高到低排；同一权重内按教材章节顺序 ----
\begin{longtable}{@{}p{0.26\linewidth}p{0.2\linewidth}p{0.4\linewidth}c@{}}
\toprule
考点 & 教材位置 & 在那里找什么 & 权重 \\
\midrule
\endhead
【考点】 & 【第〇章〇.〇节，p.〇–〇】 & 【公式（〇.〇）；例〇.〇；表〇.〇】 & 【高】 \\
\bottomrule
\end{longtable}

\section*{附录里的表（考试可能要查）}
\begin{itemize}
\item 【表名】：【p.〇】
\end{itemize}

\end{document}
